\section{Pytorch}

PyTorch è un framework di deep learning open source basato su software usato per creare reti neurali.
Combina la libreria di machine learning (ML) di Torch con un'API di alto livello basata su Python.
La sua flessibilità e facilità d'uso, tra gli altri benefici, lo hanno reso il principale framework di ML per le comunità accademiche e di ricerca.

\subsection{Tensori}

In qualsiasi algoritmo di machine learning, anche quelli applicati a informazioni apparentemente non numeriche come suoni o immagini, i dati devono essere rappresentati numericamente.
In PyTorch ciò viene ottenuto tramite tensori che fungono da unità fondamentali di dati utilizzate per il calcolo sulla piattaforma. \\

\noindent
Dal punto di vista linguistico, “tensore” è un termine generico che si riferisce ad alcune entità matematiche più familiari come :
\begin{itemize}
    \item Un vettore è un tensore monodimensionale contenente più scalari dello stesso tipo. Una tupla è un tensore monodimensionale contenente diversi tipi di dati.
    \item Una matrice è un tensore bidimensionale contenente più vettori dello stesso tipo.
    \item I tensori con tre o più dimensioni, come i tensori tridimensionali usati per rappresentare le immagini RGB negli algoritmi di computer vision, sono collettivamente denominati tensori N-dimensionali.
\end{itemize}

\noindent
Oltre a codificare gli input e gli output di un modello, i tensori di PyTorch codificano i parametri del modello: pesi e gradienti che vengono "appresi" nella machine learning.

\subsection{Grafici}
I grafici di calcolo dinamico (DCG) sono il modo in cui i modelli di deep learning vengono rappresentati in PyTorch.
Nel contesto del deep learning, questi grafici traducono essenzialmente il codice di una rete neurale in un diagramma di flusso che indica le operazioni eseguite in ciascun nodo e le dipendenze tra diversi strati nella rete, cioè la disposizione di fasi e sequenze che trasformano i dati di input in dati di output. \cite{pytorch}
